\clearpage\thispagestyle{empty}
\pdfbookmark[0]{Quellen zur Besprechung von Schulzes Schrift über das Weltsystem}{schulze1811-sources}
\begin{center}\textbf{Quellen zur Besprechung von Schulzes Schrift über das Weltsystem}\end{center}
\noindent Diese KI-übersetzte Fassung in modernem Deutsch folgt Carl Friedrich Gauß, \textit{Werke}, Band~VI (Göttingen, 1874), S.~533. Die Besprechung erschien zuerst ohne Verfassernamen unter Leipzig in den \textit{Göttingischen gelehrten Anzeigen}, 23. Februar 1811, S.~305--306. Die Werkausgabe nennt Stück~21, der Erstdruck Stück~31.
\par\smallskip\noindent Besprochen wird Gottlob Leberecht Schulzes \textit{Darstellung des Weltsystems. Ein Leitfaden für den Unterricht in der Astronomie auf Schulen} (Leipzig, Baumgärtner, 1811). Der Katalogeintrag~14/138 verzeichnet die Oktavausgabe mit XXIV Seiten Vorspann, 390 nummerierten Seiten und vier gefalteten Tafeln.
\par\smallskip\noindent \href{https://gdz.sub.uni-goettingen.de/id/PPN236007467}{Text der Werkausgabe}. Erstdruck, \href{https://images.p.sub.uni-goettingen.de/iiif/image/gdz:PPN319721507_1811:00000337/full/1650,/0/default.jpg}{S.~305} und \href{https://images.p.sub.uni-goettingen.de/iiif/image/gdz:PPN319721507_1811:00000338/full/1650,/0/default.jpg}{S.~306}. \href{https://d-nb.info/975382632/34}{Katalogeintrag zu Schulzes Buch}.
\par\smallskip\noindent Übersetzung und Satz mit KI erstellt.
\clearpage
